```latex
\documentclass[11pt]{article}

\usepackage{amsmath,amssymb,amsthm}
\usepackage{geometry}
\geometry{margin=1in}

\newtheorem{theorem}{Theorem}[section]
\newtheorem{proposition}[theorem]{Proposition}
\newtheorem{lemma}[theorem]{Lemma}
\theoremstyle{definition}
\newtheorem{definition}[theorem]{Definition}
\theoremstyle{remark}
\newtheorem{remark}[theorem]{Remark}

\newcommand{\R}{\mathbb{R}}
\newcommand{\C}{\mathbb{C}}
\newcommand{\pr}{\mathrm{pr}}

\title{A negative answer for symplectic normal crossings divisors}
\author{}
\date{}

\begin{document}

\maketitle

\begin{abstract}
We show that the answer is negative. There is a closed symplectic
$4$-manifold with a symplectic simple normal crossings divisor consisting of
two transversely intersecting symplectic surfaces whose oriented intersection
is negative. No almost complex structure tamed by the symplectic form can make
both surfaces pseudoholomorphic at their intersection. Moreover, the sign of
this crossing is invariant under any homotopy of symplectic forms for which
the two surfaces remain symplectic, so the obstruction persists after the
deformation allowed in the question.
\end{abstract}

\section{The oriented crossing obstruction}

We use the standard meaning of a symplectic simple normal crossings divisor:
it is a finite union of embedded symplectic hypersurfaces intersecting
transversely, with all intersections symplectic. Positive oriented
intersections are not part of this definition.

\begin{lemma}\label{lem:tame-orientation}
Let $(V,\omega)$ be a symplectic vector space and let $J$ be an
$\omega$-tame complex structure on $V$.
\begin{enumerate}
\item The orientation induced by $\omega^{\dim V/2}$ agrees with the complex
orientation induced by $J$.
\item If $L\subset V$ is a $J$-invariant real $2$-plane, then the
orientation induced by $\omega|_L$ agrees with the $J$-complex orientation
of $L$.
\end{enumerate}
\end{lemma}

\begin{proof}
Choose a $J$-linear identification $V\cong \C^n$ and let $\omega_0$ be the
standard constant symplectic form compatible with $J$. Define the positive
definite inner product
\[
g_0(u,v)=\omega_0(u,Jv).
\]
There is a unique $g_0$-skew-adjoint endomorphism $A$ such that
\[
\omega(u,v)=g_0(Au,v).
\]
Tameness of $J$ gives, for every $v\neq 0$,
\[
g_0(-JAv,v)=g_0(Av,Jv)=\omega(v,Jv)>0.
\]
Set
\[
A_t=(1-t)A+tJ,\qquad
\omega_t(u,v)=g_0(A_tu,v).
\]
Then
\[
g_0(-JA_tv,v)
=(1-t)g_0(-JAv,v)+t\,g_0(v,v)>0
\]
for every $v\neq0$. In particular $A_t$ is invertible, so $\omega_t$ is
symplectic for every $t$. The sign of $\omega_t^n$ is therefore constant.
At $t=1$, $\omega_1=\omega_0$ induces the $J$-complex orientation. Hence
$\omega=\omega_0$ also induces the $J$-complex orientation.

For the second statement, if $v\in L\setminus\{0\}$, then $(v,Jv)$ is a
complex-oriented basis of $L$, and
\[
\omega(v,Jv)>0.
\]
Thus the symplectic orientation of $L$ is its $J$-complex orientation.
\end{proof}

\begin{lemma}\label{lem:positive-crossing}
Let $(V,\omega)$ be a four-dimensional symplectic vector space. Suppose that
$L_1,L_2\subset V$ are transverse $J$-invariant $2$-planes for an
$\omega$-tame complex structure $J$. Then the orientation of
\[
L_1\oplus L_2,
\]
where each $L_i$ is oriented by $\omega|_{L_i}$, agrees with the orientation
of $V$ induced by $\omega^2$.
\end{lemma}

\begin{proof}
By Lemma~\ref{lem:tame-orientation}, the symplectic orientation of each
$L_i$ agrees with its $J$-complex orientation. Hence the product orientation
on $L_1\oplus L_2$ is the $J$-complex orientation of $V$. By the first part
of Lemma~\ref{lem:tame-orientation}, this agrees with the orientation
induced by $\omega^2$.
\end{proof}

\section{Invariance of the crossing sign}

Let $(M,\omega)$ be a symplectic $4$-manifold, and let $D_1,D_2$ be two
transverse symplectic surfaces meeting at $x$. Orient each $T_xD_i$ by
$\omega|_{T_xD_i}$ and orient $T_xM$ by $\omega^2$. Define the oriented
crossing sign
\[
\epsilon_x(D_1,D_2;\omega)\in\{\pm1\}
\]
to be $+1$ if the direct sum orientation on
\[
T_xD_1\oplus T_xD_2=T_xM
\]
agrees with the $\omega^2$-orientation, and $-1$ otherwise.

\begin{lemma}\label{lem:homotopy-invariance}
Let $\omega_t$, $0\le t\le1$, be a smooth path of symplectic forms on $M$.
Suppose that the fixed surfaces $D_1$ and $D_2$ are transverse and are
$\omega_t$-symplectic for every $t$. Then
\[
\epsilon_x(D_1,D_2;\omega_t)
\]
is independent of $t$.
\end{lemma}

\begin{proof}
Because $\omega_t|_{T_xD_i}$ is nondegenerate for every $t$, its induced
orientation on the fixed oriented-plane set $T_xD_i$ cannot change along the
connected interval $[0,1]$. Similarly, $\omega_t^2$ is nowhere zero, so its
induced orientation on $T_xM$ cannot change. Therefore the comparison
between the direct sum orientation and the ambient orientation is constant.
\end{proof}

\section{An explicit negative crossing}

Let $S^2$ be the unit sphere with its standard oriented area form $\sigma$.
Let
\[
f\colon S^2\longrightarrow S^2,\qquad f(x)=-x
\]
be the antipodal map. It reverses orientation, and in fact
\[
f^*\sigma=-\sigma.
\]
Choose $0<a<1$ and equip
\[
M=S^2\times S^2
\]
with the product symplectic form
\[
\omega=\pr_1^*\sigma+a\,\pr_2^*\sigma.
\]

Fix $q\in S^2$, put $p=-q$, and define
\[
D_1=S^2\times\{q\},
\qquad
D_2=\Gamma_f=\{(x,f(x)):x\in S^2\}.
\]

\begin{proposition}\label{prop:example}
The divisor $D=D_1\cup D_2$ is a symplectic simple normal crossings divisor.
Its two branches meet transversely at exactly one point, and at that point
\[
\epsilon_x(D_1,D_2;\omega)=-1.
\]
\end{proposition}

\begin{proof}
The restriction of $\omega$ to $D_1$ is $\sigma$, so $D_1$ is symplectic.
For $D_2$, let
\[
F(x)=(x,f(x)).
\]
Then
\[
F^*\omega=\sigma+a f^*\sigma=(1-a)\sigma,
\]
which is an area form because $0<a<1$. Thus $D_2$ is symplectic.

The two surfaces intersect only when $f(x)=q$. Since $f$ is a diffeomorphism,
this happens at the unique point
\[
x=(p,q).
\]
At this point,
\[
T_xD_1=T_pS^2\oplus 0,
\qquad
T_xD_2=\{(v,df_p(v)):v\in T_pS^2\}.
\]
Because $df_p$ is invertible, these planes intersect only in $0$, so the
crossing is transverse. It is therefore an ordinary simple normal crossings
crossing.

Choose a $\sigma$-positive basis $(e_1,e_2)$ of $T_pS^2$, and write
\[
\bar e_i=df_p(e_i).
\]
Since $f$ reverses orientation,
\[
\sigma_q(\bar e_1,\bar e_2)<0.
\]
A symplectically positive basis of $T_xD_1$ is
\[
B_1=\bigl((e_1,0),(e_2,0)\bigr),
\]
and a symplectically positive basis of $T_xD_2$ is
\[
B_2=\bigl((e_1,\bar e_1),(e_2,\bar e_2)\bigr).
\]
Replacing $(e_i,\bar e_i)$ by $(e_i,\bar e_i)-(e_i,0)=(0,\bar e_i)$ does not
change the orientation of the concatenated basis. Hence the orientation of
$B_1\cup B_2$ equals that of
\[
\bigl((e_1,0),(e_2,0),(0,\bar e_1),(0,\bar e_2)\bigr).
\]
Since $(\bar e_1,\bar e_2)$ is negative for the orientation of the second
sphere, this basis is negative for the product orientation
$\sigma\wedge\sigma$. Finally,
\[
\omega^2=2a\,\pr_1^*\sigma\wedge\pr_2^*\sigma
\]
induces precisely the product orientation. Therefore
\[
\epsilon_x(D_1,D_2;\omega)=-1.
\]
\end{proof}

\section{Conclusion}

\begin{theorem}\label{thm:negative-answer}
For the example $(M,D,\omega)$ above, there is no symplectic form $\omega'$
connected to $\omega$ through a path of cohomologous symplectic forms along
which $D$ remains a symplectic simple normal crossings divisor, and no
$\omega'$-tame almost complex structure $J$, such that near the crossing
point there are $J$-holomorphic local coordinates $(z_1,z_2)$ with
\[
D=\{z_1z_2=0\}.
\]
\end{theorem}

\begin{proof}
Suppose such a path $\omega_t$ and such an $\omega'=\omega_1$-tame almost
complex structure $J$ existed. At the unique crossing point $x$, the two
branches would be $\{z_1=0\}$ and $\{z_2=0\}$, possibly after interchanging
the coordinates. Because $z_i$ is $J$-holomorphic,
\[
dz_i\circ J=i\,dz_i.
\]
Consequently
\[
T_x\{z_i=0\}=\ker dz_i(x)
\]
is $J$-invariant. Thus both $T_xD_1$ and $T_xD_2$ are $J$-invariant
$2$-planes. Lemma~\ref{lem:positive-crossing} then gives
\[
\epsilon_x(D_1,D_2;\omega')=+1.
\]
Interchanging the two branches does not change the sign, since each branch
has real dimension $2$.

On the other hand, $D_1$ and $D_2$ remain $\omega_t$-symplectic for every
$t$. By Proposition~\ref{prop:example},
\[
\epsilon_x(D_1,D_2;\omega_0)=-1,
\]
and Lemma~\ref{lem:homotopy-invariance} therefore gives
\[
\epsilon_x(D_1,D_2;\omega')=-1.
\]
This is a contradiction.
\end{proof}

Thus the answer to the question is \textbf{no}. The cohomological condition
on the deformation plays no role in the argument: the crossing sign is
already invariant under any homotopy of symplectic forms for which the two
branches remain symplectic.

\end{document}
```